\begin{annexitem}{annex:p01.ec.construction}{thm:p01-ec-construction}{Finite zig-zags, leastness and intersection}
\begin{proof}
All constructions are on the fixed set \(X\). Converse is a subset
of \(X\times X\), as is its union with \(R\). The path relation
\(E\) is defined by separation from \(X\times X\). A finite
sequence has a graph contained in \(\N\times X\), so the witnesses
used in its definition can be sought within the set
\(\mathcal P(\N\times X)\). Neither a proper class of paths nor
a choice of one path for every pair is required.

\emph{Reflexivity.} For each \(x\in X\), choose length zero and
the single vertex \(x_0=x\). Its graph is the singleton
\(\{\langle0,x\rangle\}\); there is no index \(i<0\) at which
an edge is required. Thus \(xEx\). This is a path with no edges,
not a vertex-free sequence.

\emph{Containment.} If \(xRy\), the two vertices \(x_0=x,x_1=y\)
give a length-one path. Therefore \(R\subseteq E\).

\emph{Symmetry.} Given a length-\(n\) witness
\(x_0,\ldots,x_n\) from \(x\) to \(y\), define the reversed
sequence \(y_i=x_{n-i}\) for \(0\leq i\leq n\).
Its endpoints are \(y_0=y\) and \(y_n=x\). At index \(i<n\),
the original edge between \(x_{n-i-1}\) and \(x_{n-i}\) is an
\(R\)-edge in one of the two directions. The reversed pair is
therefore still an allowed zig-zag step. This argument also covers
\(n=0\), when no edge condition is tested. Hence \(xEy\)
implies \(yEx\).

\emph{Transitivity.} Suppose \(x_0,\ldots,x_m\) witnesses
\(xEy\) and \(y_0,\ldots,y_n\) witnesses \(yEz\).
Define \(u_0,\ldots,u_{m+n}\) by
\[
 u_i=x_i\quad(0\leq i\leq m),\qquad
 u_{m+j}=y_j\quad(0\leq j\leq n).
\]
At their one overlapping index the clauses agree because
\(x_m=y=y_0\). The endpoints are \(x\) and \(z\).
For \(i<m\) the adjacent pair is supplied by the first path;
for \(m\leq i<m+n\), write \(i=m+j\) with \(j<n\)
and use the second path's edge between \(y_j\) and \(y_{j+1}\).
These cases cover every required step. If either original length
is zero, the corresponding edge case is simply absent. Thus the
concatenation is a valid path and \(xEz\).

\emph{Minimality without presupposed leastness.} Let \(F\) be any
equivalence relation on \(X\) containing \(R\). For an allowed
path edge, either \(x_iRx_{i+1}\) gives \(x_iFx_{i+1}\) by
containment, or \(x_{i+1}Rx_i\) gives it by containment and
symmetry of \(F\). Induction on path length now shows
\(x_0Fx_n\). At length zero this is reflexivity of \(F\).
For the next step, the induction hypothesis gives \(x_0Fx_n\),
the final allowed edge gives \(x_nFx_{n+1}\), and transitivity
gives \(x_0Fx_{n+1}\). Hence every pair in \(E\) belongs to
\(F\). We have established \(E\subseteq F\) from the explicit
path witnesses, not from a definition that already assumes a least
relation. Since \(E\) itself is an equivalence containing \(R\),
it is least. Two least relations contain one another and are equal.

\emph{The intersection construction.} The family \(\mathcal F\)
is a subset of the set \(\mathcal P(X\times X)\). It is nonempty:
\(X\times X\) is reflexive, symmetric and transitive on \(X\)
and contains \(R\), even when \(X\) is empty.
The bounded relation \(I\) is reflexive because every \(F\in\mathcal F\)
contains each diagonal pair. If \(xIy\), symmetry in every \(F\)
gives \(yIx\). If \(xIy\) and \(yIz\), transitivity in every
\(F\) gives \(xIz\). Every \(R\)-pair belongs to every \(F\),
so \(R\subseteq I\). Thus this intersection is also an equivalence
containing \(R\). The path minimality argument yields
\(E\subseteq F\) for each \(F\in\mathcal F\), hence
\(E\subseteq I\). Conversely \(E\in\mathcal F\), so membership
in \(I\) entails membership in \(E\). Therefore \(E=I\).

\emph{Identification with reflexive-transitive closure.} By containment
and symmetry, \(E\) contains \(R\cup R^{-1}\); it is already
reflexive and transitive. If \(T\) is any reflexive and transitive
relation containing \(R\cup R^{-1}\), every edge of an \(E\)-path
belongs to \(T\). The same length induction, now using only
reflexivity and transitivity of \(T\), gives \(E\subseteq T\).
Thus \(E\) is also the least reflexive-transitive relation containing
\(R\cup R^{-1}\), which establishes its existence and its equality
with \((R\cup R^{-1})^*\) under the earlier definition of star.

\emph{Boundaries.} If \(X=\varnothing\), all three relations are
subsets of the empty product, and hence are empty. If \(R=\varnothing\),
no positive-length path can have its first allowed edge. The only
witnesses are length-zero paths, so \(E\) is exactly the diagonal
relation. Finally, if \(R\) is already an equivalence, minimality
with \(F=R\) gives \(E\subseteq R\), while containment gives
the opposite inclusion. No use of choice occurs in these cases.
\end{proof}
\end{annexitem}

\begin{annexitem}{annex:p01.ec.classes}{thm:p01-ec-classes}{Classes, their equality and the projection}
\begin{proof}
For every \(a\in A\), separation from \(A\) supplies the class
\([a]_E\). Reflexivity gives \(a\in[a]_E\), so each represented
class is nonempty. The quotient exists as the stated subset of
\(\mathcal P(A)\); it is not a collection beyond a set bound.

Suppose \(aEa'\). If \(x\in[a]_E\), then \(xEa\), and
transitivity gives \(xEa'\), hence \(x\in[a']_E\).
Symmetry gives \(a'Ea\), and the same argument yields the reverse
inclusion. Therefore \([a]_E=[a']_E\). Conversely, if the two
classes are equal, \(a\in[a]_E\) implies \(a\in[a']_E\),
which by definition means \(aEa'\).

Equal classes have a nonempty intersection, because they contain
\(a\). Conversely, if \(z\) belongs to the intersection, then
\(zEa\) and \(zEa'\). Symmetry gives \(aEz\), and transitivity
gives \(aEa'\). The established implication then makes the classes
equal. Thus any two distinct classes are disjoint. Each class is
contained in \(A\), and every \(a\in A\) belongs to its own
class, so the classes cover exactly \(A\).

The graph
\[
 \{\langle a,C\rangle\in A\times(A/E):C=[a]_E\}
\]
is a set by separation and is single-valued because the formula
\([a]_E\) specifies a unique set for each \(a\). That set belongs
to \(A/E\), so the graph is a total function \(q:A\to A/E\).
For any particular \(C\in A/E\), the quotient definition supplies
some \(a\in A\) with \(C=[a]_E=q(a)\). This proves
surjectivity without defining a simultaneous representative selector.
The equivalence \(q(a)=q(a')\) iff \(aEa'\) is exactly the class
equality already proved.

If \(A=\varnothing\), the existential condition defining \(A/E\)
has no witness, so \(A/E=\varnothing\), not \(\{\varnothing\}\).
The graph of \(q\) is the empty graph, a total surjection between
empty sets. The union of the empty family of classes is empty, and
the nonempty-class assertion has no counterinstance because there
is no class in that family. Hence the partition statement covers
the empty carrier as well.
\end{proof}
\end{annexitem}

\begin{annexitem}{annex:p01.ec.factorization}{thm:p01-ec-factorization}{The unique quotient factor without a choice of representatives}
\begin{proof}
First suppose \(f=\bar f\circ q\) for a function
\(\bar f:A/E\to B\). If \(aEa'\), the class theorem gives
\(q(a)=q(a')\). Consequently
\[
 f(a)=\bar f(q(a))=\bar f(q(a'))=f(a').
\]
Thus constancy on equivalent inputs is necessary.

Conversely suppose this constancy condition holds, and form the
displayed graph \(\bar f\) by separation from \((A/E)\times B\).
For a specified \(C\in A/E\), there exists \(c\in A\) such
that \(C=[c]_E\). Since \(c\in C\), the graph contains
\(\langle C,f(c)\rangle\), and every class has an output.
If \(\langle C,b\rangle\) and \(\langle C,b'\rangle\)
belong to the graph, let \(a,a'\in C\) witness these two
memberships. Then \(aEc\) and \(a'Ec\); symmetry and transitivity
give \(aEa'\). The hypothesis yields
\(b=f(a)=f(a')=b'\). This proves single-valuedness, so the graph
is a total function to \(B\). It also proves independence from
whichever representative might be used to describe its value.

For every \(a\in A\), membership \(a\in[a]_E\) puts
\(\langle[a]_E,f(a)\rangle\) in that graph. Hence
\(\bar f(q(a))=f(a)\), proving the factorization equation.

For uniqueness, let \(h:A/E\to B\) satisfy \(h\circ q=f\).
For each specified \(C\in A/E\), take the existential witness
\(c\) for \(C=[c]_E\). Then
\[
 h(C)=h(q(c))=f(c)=\bar f(q(c))=\bar f(C).
\]
Equality at every input gives equality of the two function graphs.
This is an argument about an arbitrary fixed class, not a choice
function assigning a witness \(c\) to all classes. The construction
of \(\bar f\) used all suitable pairs in a bounded set and required
only existence and uniqueness of each output.

If \(A=\varnothing\), its quotient is empty and this graph is
the unique empty function \(A/E\to B\). The constancy condition
and factorization equation are vacuous, and uniqueness still holds,
including when \(B=\varnothing\). For a nonempty \(A\) and empty
\(B\), no function \(f:A\to B\) satisfies the theorem's initial
hypothesis; no nonexistent input function is assumed.

For the generated-equivalence clause take \(A=X\). If a function
is constant on \(E\), then it is constant on \(R\subseteq E\).
In the other direction suppose \(xRy\) implies \(f(x)=f(y)\).
Along any zig-zag, a forward \(R\)-edge gives equal values
immediately; a backward \(R\)-edge gives them by symmetry of
equality. Induction over the finite path concatenates these equalities.
At length zero the endpoint values are identical. At the next step,
\(f(x_0)=f(x_n)\) and \(f(x_n)=f(x_{n+1})\) give the new
endpoint equality. Thus every \(xEy\) implies \(f(x)=f(y)\),
and the already proved quotient factorization applies.
\end{proof}
\end{annexitem}

\begin{annexitem}{annex:p01.ec.finite-example}{thm:p01-ec-finite-example}{Compute the finite closure and test the quotient formulas}
\begin{proof}
The displayed subsets \(C_0,C_1\) are nonempty, disjoint and
cover \(X\). Let \(T=(C_0\times C_0)\cup(C_1\times C_1)\).
This is an equivalence: every element is in a block with itself;
reversing a within-block pair stays in that block; and two
within-block pairs with a common middle element must use the same
block, since the blocks are disjoint, so their composite pair stays
in it. Every listed generator belongs to one of these products.
Thus \(R\subseteq T\), and closure minimality gives \(E\subseteq T\).

For the reverse inclusion, the generators give \(0E1\), \(2E1\)
and \(2E3\). Symmetry and transitivity give \(3E1\), while
reflexivity gives \(1E1\). Every \(a\in C_0\) therefore satisfies
\(aE1\). For any \(a,b\in C_0\), symmetry of \(bE1\) and
transitivity with \(aE1\) give \(aEb\). Similarly \(4E4\)
and the reversed generator \(5E4\) show that any two elements of
\(C_1\) are equivalent. Hence \(T\subseteq E\), proving \(E=T\).

There are \(4\cdot4=16\) pairs in \(C_0\times C_0\) and
\(2\cdot2=4\) in \(C_1\times C_1\). The products are disjoint,
so \(E\) contains twenty ordered pairs. Its classes at inputs
\(0,1,2,3\) are exactly \(C_0\), because the only related
elements lie in that block; its classes at \(4,5\) are exactly
\(C_1\). This proves the displayed quotient and all six projection
values, including \([0]_E=[1]_E\) and \(0E3\).

The relation consisting of the diagonal on \(X\) together with
\(R\) is reflexive and transitive. For transitivity, a diagonal
leg changes no endpoint, while two successive nondiagonal \(R\)-legs
are impossible: the set of sources \(\{0,2,4\}\) is disjoint
from the set of targets \(\{1,3,5\}\). Any reflexive relation
containing \(R\) must contain this union. It is therefore the
least reflexive-transitive relation containing \(R\), hence is
\(R^*\). The pair \(\langle0,3\rangle\) belongs neither to
the diagonal nor to \(R\), so \(0R^*3\) is false. In contrast,
the sequence \(0,1,2,3\) is an allowed zig-zag: its middle
step uses \(2R1\) backwards.

For each generator the two values of \(f\) agree: both are
\(u\) on the three \(C_0\) generators and both are \(v\)
on the \(C_1\) generator. The generated-equivalence theorem
therefore supplies the unique factor. Its graph is explicitly
\(\{\langle C_0,u\rangle,\langle C_1,v\rangle\}\),
which verifies both stated values and the equation
\(f=\bar f\circ q\) at every element of \(X\).

The function \(g:X\to X\) with \(g(x)=x\) has
\(g(0)=0\ne1=g(1)\), although \(0E1\). It violates
the necessary constancy condition. Directly, the proposed
formula would assign \(\bar g(C_0)=0\) when using representative
zero and \(\bar g(C_0)=1\) when using representative one.
One input cannot have both distinct outputs; the formula is not a
function on the quotient.

The separately specified rule \(t(C)=\min C\) is well-defined
on this two-element quotient: inspecting the two finite subsets
of \(\N\) gives least elements zero and four respectively.
Its graph is \(\{\langle C_0,0\rangle,\langle C_1,4\rangle\}\).
Then \(q(t(C_0))=q(0)=C_0\) and
\(q(t(C_1))=q(4)=C_1\), so \(q\circ t\) is the quotient's
identity function. But \((t\circ q)(1)=t(C_0)=0\ne1=g(1)\).
Thus this valid finite representative selector is not a factor
of \(g\), and neither its existence nor a preferred choice of
representative supplies the missing independence argument.
\end{proof}
\end{annexitem}
